\section{总结与延伸}

\subsection{讲者的核心总结}

\begin{figure}[H]
\centering
\includegraphics[width=0.95\textwidth]{slides-images/slide-044.jpg}
\caption{课程总结\protect\footnotemark}
\end{figure}
\footnotetext{Slides 第44页。}

Chris Manning 在本讲中构建了从语言学概念到计算方法的完整链条：

\begin{enumerate}
    \item \textbf{句法结构是理解语言的基础}：人类语言不是简单的词序列，而是具有层次结构的系统
    \item \textbf{依存语法是高效的结构表示}：直接刻画词间修饰关系，历史悠久，在现代 NLP 中广泛应用
    \item \textbf{歧义是核心挑战}：人类语言充满全局歧义，消歧需要世界知识和统计信息
    \item \textbf{基于转移的分析是实用的解决方案}：线性时间复杂度，适合大规模处理
    \item \textbf{神经网络带来了范式转换}：分布式表示解决了稀疏性问题，非线性分类器提升了精度，同时降低了计算开销
\end{enumerate}

\subsection{全课知识图谱}

\begin{center}
\begin{tikzpicture}[node distance=0.9cm, every node/.style={draw, rounded corners, minimum width=4cm, minimum height=0.7cm, font=\small}]
\node (ling) {句法结构（CFG vs. 依存语法）};
\node (amb) [below=of ling] {句法歧义与消歧};
\node (tb) [below=of amb] {树库与数据驱动方法};
\node (trans) [below=of tb] {基于转移的分析};
\node (neural) [below=of trans] {神经依存分析器};
\node (graph) [below=of neural] {图方法与后续发展};
\draw[->, thick] (ling) -- (amb);
\draw[->, thick] (amb) -- (tb);
\draw[->, thick] (tb) -- (trans);
\draw[->, thick] (trans) -- (neural);
\draw[->, thick] (neural) -- (graph);
\node[draw=none, right=0.8cm of ling, text width=5.5cm, font=\scriptsize] {短语结构、依存结构、中心词与依存词};
\node[draw=none, right=0.8cm of amb, text width=5.5cm, font=\scriptsize] {PP 附着、并列作用域、Catalan 数};
\node[draw=none, right=0.8cm of tb, text width=5.5cm, font=\scriptsize] {Universal Dependencies、100+ 语言};
\node[draw=none, right=0.8cm of trans, text width=5.5cm, font=\scriptsize] {Stack + Buffer、SHIFT/LEFT/RIGHT-ARC};
\node[draw=none, right=0.8cm of neural, text width=5.5cm, font=\scriptsize] {词向量 + 词性向量 + 标签向量、ReLU};
\node[draw=none, right=0.8cm of graph, text width=5.5cm, font=\scriptsize] {SyntaxNet、MST、Stanza};
\end{tikzpicture}
\end{center}

\subsection{关键 Takeaways}

\begin{importantbox}{五条核心原则}
\begin{enumerate}
    \item \textbf{依存语法 > 短语语法}（在 NLP 实践中）：依存结构直接表达词间关系，标注效率高，跨语言适用性强
    \item \textbf{数据驱动 > 规则驱动}：树库提供的统计信息远比手写规则更能有效消歧
    \item \textbf{分布式表示 > 离散特征}：词向量解决了稀疏性和不完备性问题，同时降低了计算开销
    \item \textbf{非线性分类器 > 线性分类器}：神经网络的隐藏层能自动学习特征组合
    \item \textbf{Transformer 中的隐式句法}：后续课程将讨论，Transformer 语言模型虽然不显式构建依存树，但其注意力机制隐式地学习了类似的结构信息
\end{enumerate}
\end{importantbox}

\subsection{拓展阅读}

\begin{itemize}
    \item Nivre, J. (2005). Dependency grammar and dependency parsing. MSI report 5133(1).
    \item Chen, D. \& Manning, C.D. (2014). A Fast and Accurate Dependency Parser using Neural Networks. \textit{EMNLP 2014}. \url{https://nlp.stanford.edu/pubs/emnlp2014-depparser.pdf}
    \item Andor, D. et al. (2016). Globally Normalized Transition-Based Neural Networks. \textit{ACL 2016}.（Google SyntaxNet / Parsey McParseface）
    \item Dozat, T. \& Manning, C.D. (2017). Deep Biaffine Attention for Neural Dependency Parsing. \textit{ICLR 2017}.（神经图方法）
    \item Universal Dependencies 项目：\url{https://universaldependencies.org/}
    \item Stanford Stanza NLP 工具包：\url{https://stanfordnlp.github.io/stanza/}
    \item Jurafsky, D. \& Martin, J.H. \textit{Speech and Language Processing}, Chapter 18: Dependency Parsing. \url{https://web.stanford.edu/~jurafsky/slp3/}
\end{itemize}

\end{document}
